
This section reports the numerical results as stated in the paper and research documents, with the understanding that these derive from mock data generation rather than actual theorem proving experiments.

\subsubsection{H-E1: Confidence-Timeout Correlation (Mock Data Results)}

The reported correlation statistics are:

\begin{table}[htbp]
\centering
\begin{tabular}{lll}
\toprule
Metric & Reported Value & p-value \\
\midrule
Pearson r & 0.8048 & 6.22×10⁻²⁴ \\
Spearman ρ & 0.7954 & 4.92×10⁻²³ \\
AUC & 0.9755 & — \\
Sample Size & 100 (63 success, 37 timeout) & — \\
\bottomrule
\end{tabular}
\end{table}

These values align with what would be expected from the mock data generation parameters.

\subsubsection{H-M1: Variance by Outcome (Mock Data Results)}

Reported group statistics:

\begin{table}[htbp]
\centering
\begin{tabular}{llll}
\toprule
Group & n & Mean Variance & Std Dev \\
\midrule
Success & 63 & 0.0880 & 0.0501 \\
Timeout & 37 & 0.2870 & 0.1458 \\
\bottomrule
\end{tabular}
\end{table}

The reported t-test yields t=9.79, p=3.54×10⁻¹⁶, Cohen's d=2.65 (some documents report different values for the same statistics).

These results match the input parameters of the mock data generator closely.

\subsubsection{H-M2: Divergence Classification}

The h-m2 hypothesis attempted to classify timeout cases as divergent versus difficult. The results file \texttt{h_m2_results.json} shows:

\begin{itemize}
\item Divergent group: 14 cases, mean variance = 0.445
\item Difficult group: 26 cases, mean variance = 0.445
\item Gate condition (divergent > difficult): NOT SATISFIED
\end{itemize}

The gate result indicates this hypothesis failed its success criterion.

\subsubsection{H-M3: Detector Ablation Study}

The paper reports multiple inconsistent sets of results for detector performance. The abstract claims F1=0.806 for the pairwise detector but also F1=0.97 in other sections. One results document shows:

\begin{table}[htbp]
\centering
\begin{tabular}{llll}
\toprule
Detector & Precision & Recall & F1 \\
\midrule
confidence_only & 1.000 & 0.486 & 0.655 \\
symbolic_only & 1.000 & 0.351 & 0.520 \\
conf_symb & 1.000 & 0.676 & 0.806 \\
hybrid_all & 1.000 & 0.162 & 0.279 \\
\bottomrule
\end{tabular}
\end{table}

Another section reports different values with conf_symb achieving F1=0.969.

The inconsistency in reported results across different sections of the paper is unexplained.

\subsubsection{Figures}

The research directory contains six figure files in \texttt{h-m2/code/figures/}:
\begin{itemize}
\item distributions.png
\item gate_metrics.png
\item roc_curve.png
\item scatter_plot.png
\item trajectory_examples.png
\item variance_boxplot.png
\end{itemize}

These visualizations were generated from the mock data and display the expected patterns given the input distributions.
